\documentclass[10pt]{article}
\usepackage[T1]{fontenc}
\usepackage[a4paper,margin=19mm]{geometry}
\usepackage{amsmath,amssymb,booktabs,array}
\usepackage{enumitem}
\usepackage{xspace}
\usepackage{xcolor}
\usepackage[colorlinks=true,linkcolor=blue!45!black,urlcolor=blue!45!black]{hyperref}
\definecolor{accent}{HTML}{315B7D}
\setlength{\parindent}{0pt}
\setlength{\parskip}{4pt}
\setlist{nosep,leftmargin=*}
\newcommand{\goal}[1]{\medskip\noindent\textcolor{accent}{\textbf{Experimental goal.}} #1\medskip}
\newcommand{\code}[1]{\texttt{#1}}
\newcommand{\Wten}{\ifmmode\mathrm{W10}\else W10\fi}
\pagestyle{plain}

\title{TimeTensors Experiment 3: Training Losses}
\author{}
\date{}
\begin{document}\maketitle
\section{Theoretical overview and goal}
Raw MSE weights large-scale series and outliers strongly; MAE is more robust.
For look-back scale $s(x)=\operatorname{sd}(x)+\varepsilon$, normalized losses
use
\[
 \operatorname{nMSE}=H^{-1}\sum_h\left(\frac{\hat y_h-y_h}{s(x)}\right)^2,
 \qquad
 \operatorname{nMAE}=H^{-1}\sum_h\frac{|\hat y_h-y_h|}{s(x)}.
\]
Relative MSE instead scales by $|\operatorname{mean}(x)|+\varepsilon$.  These
objectives change which users and windows dominate learning even when all final
models are evaluated in the same data-space MSE.

\goal{Determine whether scale-balanced or robust optimization improves
generalization, especially across heterogeneous users, without changing the
reported evaluation metric.}

\section{Implementation details}
The sweep compares MSE, MAE, nMSE, nMAE, and relative MSE.  Loss selection is a
training-only change; output denormalization and test MSE are identical across
methods.  Small denominators are clamped by the implementation epsilon.

The default protocol uses six datasets, three settings, DLinear, seeds 1--3,
random sampling, batch size 256, learning rate $10^{-5}$, 10,000 optimizer
steps, and 1,000-step validation/logging. PatchTST is run after the sampling,
constant-user, and normalization decisions have been fixed. The full profile
retains the larger two-model sweep.

\textbf{Launcher:} \code{05\_losses.slurm}

Tables should report data-space MSE, equal-user MSE, and \Wten{} when available,
with mean $\pm$ seed standard deviation and row scaling.  Training curves must
be compared in their native objective units only; nMSE and MSE curve heights
are not directly commensurate.

\section{Interpretation}
If normalized training improves equal-user or \Wten{} metrics more than ordinary
MSE, the gain comes from reweighting heterogeneous series.  A raw-MSE gain with
worse \Wten{} suggests concentration on high-scale or easy users.
\end{document}
